\begin{abstract}
Post-quantum communication in IoT requires secure session recovery after restarts. Challenges include (1) restoring obsolete cryptographic state, (2) losing resumption context, and (3) duplicating effects after lost acknowledgements. In this work, we evaluate recovery combining authenticated checkpoints, fresh traffic keys, and persistent deduplication. Our approach resumes sessions or uses full authentication and saves effects before acknowledgement. We evaluate this through fault injection and case studies in Mininet-WiFi with classical, hybrid, and post-quantum TLS 1.3. With gateway restarts and acknowledgement suppression at every eligible reading and 0\% configured packet loss, the post-quantum configuration saves 67.87\% of IPv4 bytes per confirmed reading with preserved context and increases costs by 14.24\% with lost context, relative to full authentication. The server chain includes the root CA certificate. Deduplication yielded zero duplicate database effects, against 2,700 without suppression. These results link recovery costs to preserved state.
\end{abstract}

\begin{resumo}
A comunicação pós-quântica em redes IoT exige recuperação segura de sessões após reinícios. Os desafios incluem (1) restaurar estado obsoleto, (2) perder o contexto de retomada e (3) duplicar efeitos após perda de confirmações. Neste trabalho, avaliamos uma abordagem que combina checkpoints autenticados, chaves de tráfego novas e deduplicação persistente. Nossa abordagem retoma sessões ou utiliza autenticação completa e salva efeitos antes da confirmação. Avaliamos a solução por injeção de falhas e estudos de caso no Mininet-WiFi com TLS 1.3 clássico, híbrido e pós-quântico. Com reinícios do gateway e supressão da confirmação em cada leitura elegível, sob perda de pacotes configurada em 0\%, a configuração pós-quântica reduziu em 67,87\% os bytes IPv4 por leitura confirmada com contexto preservado e aumentou o custo em 14,24\% com contexto perdido, frente à autenticação completa. A cadeia do servidor inclui o certificado da CA raiz. A deduplicação resultou em zero efeitos duplicados no banco, contra 2.700 sem supressão. Os resultados relacionam o custo da recuperação ao estado preservado.
\end{resumo}
